% =============================================================
%  chapter2_abstract.tex
%  Shared abstract text for Chapter 2 (goethite hydrothermal step
%  heat experiment). Kept separate from any \begin{abstract}/
%  \section{Abstract} wrapper for the same reason as Chapter 1's
%  chapter1_abstract.tex: the surrounding heading/environment is
%  not class-agnostic, but this raw text is.
%
%  \input{} by BOTH:
%    - Chapter2/Writing/chapter2goethite.tex  (standalone draft)
%    - Chapter2/chapter2.tex                     (dissertation chapter)
%  Edit this file, not the two wrappers, to change the abstract text.
% =============================================================
U-Th/He geochronology of goethite (FeOOH) is a powerful method to
quantitatively date weathering and erosion processes on Earth's surface;
however, precise interpretation of goethite U-Th/He dates is challenging
because uncertainty persists regarding the mineralogical and environmental
conditions under which radiogenic helium is retained and how it may be lost over
time. Standard vacuum step heating experiments have been successfully
used on a variety of mineral phases to characterize helium migration,
but the results of such experiments on goethite have been dominated by
dehydration-induced helium loss, raising questions about their relevance to
helium retention and loss in nature \citep{farley2024}. Observations
that thermal decomposition of goethite is greatly slowed in the presence
of high water pressure led us to develop a new method to conduct helium
step heat experiments under hydrothermal conditions (up to 86 bars and
300°C) within standard experimental setups. Using this method, we report
results from the first helium step heat experiment of proton-irradiated
goethite under hydrothermal conditions. Irradiated goethite sample
Win-06-01b of the Flinders Ranges, South Australia and distilled
water were encapsulated in a pure silica glass ampoule and were
subjected to multiple prograde and retrograde thermal cycles over a
temperature range of 170°C to 250°C under water vapor saturated
conditions. Step durations varied from 2 hours to over a week, with
low-temperature (\textless120°C) permeation steps bracketing
higher-temperature steps to account for residual helium within the
ampoule. Because 4He yields from those permeation steps were at blank
level, permeation lag was removed by the bounding method, which
constrains the ampoule's residual helium from above using the rate
constant measured during the experiment itself; the permeation step
method, applied to 3He alone, reproduces the same result. The results show apparent helium diffusivities approximately
seven log units lower than those observed in vacuum, with apparent
diffusivities following a consistent linear array on an Arrhenius plot
for cumulative fractional 3He loss exceeding 7\%. Delta plots generated
using this reference array with a calculated activation energy of \ChTwoEa
± \ChTwoEaErr kJ/mol are consistent with a volume-diffusion helium loss
mechanism with MDD (Multiple Diffusion Domain) systematics. These
results show high helium retentivity in natural goethite under
hydrothermal conditions that extrapolate to robust helium retention at
surface conditions even over billion-year timescales. The method
presented in this work can be systematically applied to experimentally
investigate effects of crystallite size, chemical composition, and
radiation damage to refine these results and to assess the applicability
of existing vacuum step heat data to helium retention and loss in
goethite in nature.
